% =============================================================================
% practica3.tex --- Chapter: Practice 3
%
% Compile standalone:
%   pdflatex -shell-escape practica3.tex   (run twice for cross-references)
%
% Compile as part of the full report:
%   pdflatex -shell-escape __memoria.tex   (from docLaTex root)
% =============================================================================

\documentclass[../../__memoria.tex]{subfiles}

\begin{document}
\standalonecover

% ── Graphics path relative to THIS file when compiled standalone ─────────────
% When compiled from __memoria.tex the global \graphicspath already covers
% {img/} via \subfix. The extra paths below only activate for standalone:
%   - ../../Imagenes/  → needed for logo_uco (used in fancyhdr header)
%   - img/             → figures for this practice
\graphicspath{{\subfix{../../Imagenes/}}{../../Imagenes/}{\subfix{img/}}{img/}}

\chapter{Cuaderno de Laboratorio --- Práctica 3: Metodología de Evaluación y Rigor Científico}

\section*{Objetivo}
Aprender a realizar evaluaciones honestas en condiciones difíciles, especialmente con datos desbalanceados,
y detectar errores metodológicos graves como la fuga de datos (\textit{data leakage}).

\section*{Material de partida}
Se parte de:
\begin{itemize}
    \item Un dataset sintético de \textbf{1000 pacientes} con \textbf{5 variables predictoras}.
    \item Una clase positiva muy minoritaria (\textbf{2\%} de los casos).
\end{itemize}

% =============================================================================
\section{Tarea 1 --- La lotería de la partición aleatoria}
% =============================================================================

\subsection*{Configuración común}
Modelo usado: \texttt{LogisticRegression} (con \texttt{class\_weight='balanced'})\\
Métricas usadas: \texttt{Nº positivos en test}, \texttt{\% clase 1}\\
Semilla(s): \texttt{0, 1, 2, 3, 4} (5 ejecuciones aleatorias)

\vspace{0.3cm}
\noindent\textbf{Hipótesis inicial (antes de ejecutar):}
Es probable que el número de positivos en test sea muy inestable.
Al haber solo un 2\% de casos positivos (20 de 1000), una partición aleatoria
de 200 ejemplos (20\%) podría capturar por puro azar una cantidad muy variable
de estos casos, o incluso ninguno, invalidando la representatividad del
conjunto de test.
\vspace{0.3cm}

\subsection{Registro de resultados (5 particiones aleatorias)}
\begin{center}
\begin{tabular}{|c|c|}
\hline
Ejecución & Positivos en test \\
\hline
1 & 4 \\
2 & 3 \\
3 & 5 \\
4 & 4 \\
5 & 3 \\
\hline
\end{tabular}
\end{center}

\begin{figure}[H]
    \centering
    \includegraphics[width=0.85\linewidth]{tarea1_loteria}
    \caption{Variabilidad del número de ejemplos positivos en el conjunto de test según la semilla aleatoria (partición sin estratificar).}
    \label{fig:t1_loteria}
\end{figure}

\noindent\textbf{Análisis:}
\begin{itemize}
    \item \textit{¿Ha cambiado mucho el número de positivos en test entre ejecuciones?}
          Sí, el número oscila entre 3 y 5. Aunque parece una diferencia pequeña en
          valor absoluto, supone una variación relativa significativa ($\pm 25\%$) en un
          conjunto de test tan pequeño.
    \item \textit{¿Sería fiable la evaluación si el azar dejara solo 0 o 1 positivos en test?}
          No. Con 0 o 1 positivos, métricas como F1-score o Recall para la clase
          minoritaria serían imposibles de calcular o carecerían de significado
          estadístico, dando resultados extremos (0 o 100\%) que no reflejan la
          capacidad real del modelo.
\end{itemize}

% =============================================================================
\section{Tarea 2 --- Partición estratificada}
% =============================================================================

\subsection{Registro por folds (\texttt{StratifiedKFold}, \texorpdfstring{$k=10$}{k=10})}
\begin{center}
\begin{tabular}{|c|c|c|c|}
\hline
Fold & Positivos en test & Tamaño del fold & Proporción clase 1 \\
\hline
1  & 2 & 100 & 2{,}0\% \\
2  & 2 & 100 & 2{,}0\% \\
3  & 2 & 100 & 2{,}0\% \\
4  & 2 & 100 & 2{,}0\% \\
5  & 2 & 100 & 2{,}0\% \\
6  & 2 & 100 & 2{,}0\% \\
7  & 2 & 100 & 2{,}0\% \\
8  & 2 & 100 & 2{,}0\% \\
9  & 2 & 100 & 2{,}0\% \\
10 & 2 & 100 & 2{,}0\% \\
\hline
\end{tabular}
\end{center}

\noindent Desviación típica de la proporción de clase 1 entre folds: \textbf{0{,}0000\%}
(prácticamente nula, lo que confirma que la estratificación funciona correctamente).

\begin{figure}[H]
    \centering
    \includegraphics[width=0.85\linewidth]{tarea2_estratificada}
    \caption{Proporción de la clase positiva por fold con \texttt{StratifiedKFold}. La línea roja indica la proporción global del dataset (2\%).}
    \label{fig:t2_estratificada}
\end{figure}

\noindent\textbf{Análisis:}
\texttt{StratifiedKFold} garantiza que cada fold mantenga la misma proporción de
clases que el dataset original (2\%). En datasets muy desbalanceados, esto evita
que algunos folds se queden sin muestras de la clase minoritaria, asegurando que
el modelo sea evaluado siempre sobre una distribución representativa y reduciendo
la varianza debida al azar.

% =============================================================================
\section{Tarea 3 --- Comparativa de la varianza}
% =============================================================================

\subsection{Resumen de resultados (10 iteraciones, 10 folds)}
\begin{center}
\begin{tabular}{|l|c|c|c|c|}
\hline
Método & Accuracy media & std(Accuracy) & F1-macro media & std(F1) \\
\hline
\texttt{KFold}           & 0{,}5771 & 0{,}0090 & 0{,}3829 & 0{,}0051 \\
\texttt{StratifiedKFold} & 0{,}5783 & \textbf{0{,}0085} & 0{,}3857 & \textbf{0{,}0050} \\
\hline
\end{tabular}
\end{center}

\begin{figure}[H]
    \centering
    \includegraphics[width=0.88\linewidth]{tarea3_varianza}
    \caption{Distribución de Accuracy y F1-macro a lo largo de 10 iteraciones de validación cruzada. \texttt{StratifiedKFold} produce cajas más estrechas (menor varianza).}
    \label{fig:t3_varianza}
\end{figure}

\noindent\textbf{Interpretación:}
\begin{itemize}
    \item \textit{¿Qué método produce métricas más estables?}
          El método \texttt{StratifiedKFold} produce métricas más estables,
          como se observa en la menor desviación típica (\textit{std})
          tanto en Accuracy como en F1-macro.
    \item \textit{¿Dónde se nota más la mejora, en Accuracy o en F1?}
          La mejora en la estabilidad es notable en ambas métricas, pero el
          rigor metodológico es especialmente crítico en F1-macro, ya que
          esta métrica penaliza el mal rendimiento en la clase minoritaria,
          que es la más sensible a particiones no representativas.
\end{itemize}

% =============================================================================
\section{Tarea 4 --- Detección de \textit{data leakage}}
% =============================================================================

\subsection{Comparativa: con y sin variable trampa}
\begin{center}
\begin{tabular}{|l|c|c|}
\hline
Escenario & Accuracy & F1-macro \\
\hline
Sin variable trampa (\textit{evaluación real})  & 0{,}5850 & 0{,}3887 \\
Con variable trampa (\texttt{ID\_Hospital\_Filtro}) & \textbf{1{,}0000} & \textbf{1{,}0000} \\
\hline
\end{tabular}
\end{center}

\begin{figure}[H]
    \centering
    \includegraphics[width=0.85\linewidth]{tarea4_leakage}
    \caption{Efecto del \textit{data leakage}. Las barras rojas corresponden al escenario con la variable trampa; las verdes, al escenario real sin ella.}
    \label{fig:t4_leakage}
\end{figure}

\noindent\textbf{Preguntas:}
\begin{itemize}
    \item \textit{¿Qué variable está filtrando información del objetivo?}
          La variable \texttt{ID\_Hospital\allowbreak\_Filtro}.
    \item \textit{¿Por qué esa variable invalida la evaluación?}
          Ha sido construida con una correlación perfecta (o casi perfecta) con
          la clase objetivo. El modelo aprende a usar este ``filtro'' en lugar de
          las variables clínicas reales, obteniendo un rendimiento del 100\% que
          es imposible en la práctica real.
    \item \textit{¿Diferencia entre un modelo ``bueno'' y uno contaminado por leakage?}
          Un modelo bueno aprende patrones generalizables a partir de predictores
          legítimos. Un modelo con \textit{leakage} hace trampa: usa información
          que no estará disponible en producción, por lo que fallará completamente
          al enfrentarse a nuevos datos.
\end{itemize}

% =============================================================================
\section{Reto --- Auditoría metodológica}
% =============================================================================

\begin{enumerate}
    \item \textbf{Peligros de la partición aleatoria:} Puede generar conjuntos de
          test vacíos de clase minoritaria, impidiendo una evaluación estadística
          significativa y dando una falsa sensación de precisión mediante el
          ``acierto por mayoría''.
    \item \textbf{Aporte de la estratificación:} Garantiza la representatividad
          de la muestra en cada fold, permitiendo comparar modelos de forma justa
          y estable.
    \item \textbf{Data leakage:} Hace inútil el experimento porque el modelo no
          aprende a resolver el problema real, sino a identificar una variable que
          ya contiene la respuesta, invalidando su uso futuro en producción.
    \item \textbf{Buenas prácticas:} Usar siempre estratificación, auditar
          variables en busca de correlaciones sospechosas, emplear métricas
          robustas (F1, AUC-ROC) y documentar la reproducibilidad mediante
          semillas y validación cruzada repetida.
\end{enumerate}

% =============================================================================
\section{Conclusión final}
% =============================================================================

La evaluación de modelos de IA no consiste en ejecutar un comando y leer un
porcentaje: requiere rigor metodológico. En problemas con datos desbalanceados,
el azar es un adversario que la estratificación ayuda a controlar.
Además, el \textit{data leakage} es un error silente y catastrófico: métricas
perfectas son casi siempre señal de un error de diseño, no de un modelo
excepcional. La honestidad en la evaluación es la base de cualquier sistema de
IA confiable y reproducible.

\end{document}
